% ---------------------------------------------------------------------------
% Section 4 -- The Framework: OUTLINE / SKETCH
%
% This file is a talking-points skeleton, not prose. It is NOT \input by
% main.tex; the written Sec.4 currently lives inline in main.tex (lines
% 524-767). Edit the bullets here to decide the shape, then write against it.
%
% Each subsection carries:
%   PURPOSE  -- the one question this subsection answers for the reader
%   POINTS   -- what to say, in order
%   OMIT     -- what belongs to another section, so it does not creep back in
%
% Revised structure:
%   4.1 Framework Overview -- static vs. dynamic and why the split exists
%   4.2 Static Layer
%       4.2.1 Tool Impact
%       4.2.2 Asset Sensitivity
%       4.2.3 Blast Radius
%   4.3 Dynamic Layer
%
% Enforcement is intentionally out of scope. The framework assesses risk and
% exposes an auditable score/band; the organization decides how to act on it.
%
% Budget: 1,400-1,700 words total (WRITING-GUIDE.md Sec.5). Suggested split in
% the comment above each subsection. Every number must appear verbatim in
% GROUND-TRUTH.md before it goes into prose.
% ---------------------------------------------------------------------------
\section{The Framework}
\label{sec:framework}

% --- Lead-in, no heading. ~50-60 words. ------------------------------------
% PURPOSE: tell the reader immediately that the framework separates capability
% from runtime context.
\begin{itemize}
  \item \framework{} computes the risk of a permitted invocation in two layers.
  \item The \emph{static layer} estimates the inherent consequence of a
        (tool, asset) pairing before any particular call occurs.
  \item The \emph{dynamic layer} refines that baseline at request time using the
        concrete arguments, session context, and the calling agent's history.
  \item The framework produces an auditable risk score/band; it does not define
        an allow/review/block enforcement policy.
\end{itemize}

% -------------------------------------------------------------------------
\subsection{Framework Overview}
\label{sec:overview}
% ~350-400 words.
% PURPOSE: explain the two-layer architecture, what each layer knows, why they
% are separated, and how they compose into rho(c).
\begin{itemize}
  \item \textbf{The central design choice: separate capability from context.}
        The same tool can have very different consequences depending on the
        asset it touches, while the same (tool, asset) pair can become more or
        less concerning depending on the concrete request and recent behavior.
        These are different questions and should not be collapsed into one
        opaque model decision.

  \item \textbf{Why the static layer exists.}
    \begin{itemize}
      \item It captures facts that are stable across calls: what operation the
            tool performs, how sensitive the asset is, and how broadly one use
            of that tool can affect it.
      \item These values can be computed ahead of time and reused across many
            invocations, giving every (tool, asset) pair a baseline consequence.
      \item This makes the baseline auditable: an operator can inspect which
            factor produced a high score rather than receiving one unexplained
            number.
    \end{itemize}

  \item \textbf{Why the dynamic layer is separate.}
    \begin{itemize}
      \item Static analysis cannot see the actual arguments, whether the call is
            unusual for this agent, or whether it is part of a risky sequence.
      \item Those signals exist only when a request arrives, so they belong to a
            request-time layer rather than to the precomputed baseline.
      \item Keeping them separate also prevents transient behavioral evidence
            from changing the underlying classification of a tool or asset.
    \end{itemize}

  \item \textbf{The two existing inputs.}
    \begin{itemize}
      \item the tool catalog the server advertises over \texttt{tools/list};
      \item the organization's policy/register for the assets exposed through
            the server, including named classification classes and recognition
            rules, but not framework-specific severity numbers.
      \item Neither input is created only for \framework{}; this is part of the
            deployability argument.
    \end{itemize}

  \item \textbf{How the layers compose.}
        The static layer computes the three-factor baseline in
        Eq.~\ref{eq:score}. The dynamic layer then adjusts that baseline using
        request-time evidence to produce $\rho(c)$.

  \item \textbf{Why the static score has three factors.}
        Tool impact is the \emph{verb} --- what the operation does; asset
        sensitivity is the \emph{noun} --- what it touches; blast radius is the
        \emph{scope} --- how far the consequence can propagate. Each captures a
        different source of consequence that the other two cannot express.

  \item \textbf{Why multiply.}
        The three 1--5 primitives are multiplied, producing a range of
        $[1,125]$. Multiplication preserves interaction between the factors: a
        call reaches the top of the scale only when it is severe across all
        three dimensions, while a low value on one axis remains influential.

  \item \textbf{Output and scope.}
        The result is mapped to an ordinal risk band. \framework{} stops at risk
        assessment: how an organization translates that assessment into an
        operational action is outside this framework.
\end{itemize}
% OMIT: detailed factor rubrics -- 4.2; detailed runtime signals -- 4.3;
%       enforcement thresholds or allow/review/deny logic -- out of scope.

% -------------------------------------------------------------------------
\subsection{The Static Layer}
\label{sec:static}
% ~80-120 words before the three subsubsections.
% PURPOSE: define exactly what the static layer produces and how often it runs.
\begin{itemize}
  \item The static layer constructs a reusable baseline for each (tool, asset)
        pairing exposed by the server.
  \item It scores three primitives on 1--5: Tool Impact $I(t)$, Asset
        Sensitivity $S(r)$, and Blast Radius $B(t,r)$.
  \item Keep Eq.~\ref{eq:score}: the static score is the product of these three
        values, before request-time evidence is considered.
  \item Conceptually: impact is primarily tool-level, sensitivity is
        asset-level, and blast radius is pair-level because propagation depends
        on both the operation and the asset it reaches.
  \item Static computation is repeated when the tool catalog, asset register, or
        classification policy changes; it is not recomputed from scratch for
        every invocation.
\end{itemize}

% -------------------------------------------------------------------------
\subsubsection{Tool Impact}
\label{sec:impact}
% ~350-400 words + Table 2.
% PURPOSE: explain how the framework estimates the inherent severity of the
% operation exposed by a tool.
\begin{itemize}
  \item \textbf{The question.} What can one invocation of this tool do to its
        target, independent of which particular asset is selected?

  \item \textbf{The ladder.} Keep Table~\ref{tab:ladder}: 1 none, 2 metadata,
        3 content read or limited write, 4 ordinary write, 5 removal or
        execution.

  \item \textbf{The vocabulary, with the counts.} 281 patterns
        (16 none / 52 metadata / 46 read / 122 write / 45 remove), plus 119
        ambiguous single words matched \emph{against the tool name only},
        10 generic read verbs, and 9 parameter detectors.

  \item \textbf{Why name-only scoping matters.} Matching ambiguous verbs across
        the whole description can make a harmless listing or descriptive tool
        appear destructive merely because its documentation mentions another
        action.

  \item \textbf{The 3/4 boundary is borrowed, not invented.} Use the analogy to
        HTTP \texttt{PATCH} vs. \texttt{PUT} and CVSS~v4.0
        \texttt{VI:L} vs. \texttt{VI:H}: a bounded modification belongs below
        an ordinary full write. Example: writing one sentence into a file is 3;
        writing/replacing the file is 4.

  \item \textbf{Two deliberate absences.}
    \begin{itemize}
      \item breadth indicators such as all/every/bulk/batch, globs, or recursion
            do not raise Tool Impact because breadth belongs to Blast Radius;
      \item parameter detectors identify capabilities whose actual use depends
            on the request arguments, so they do not by themselves move the
            static impact tier.
    \end{itemize}

  \item \textbf{MCP annotations are supporting evidence, not a bound.}
        An annotation such as \texttt{destructiveHint} may corroborate a rule,
        but it neither sets nor caps the impact score.

  \item \textbf{Rules-to-model hand-off.} Fold the former separate hand-off
        subsection here:
    \begin{itemize}
      \item confidence 0.35 when no tier verb matched, 0.8 when a verb fired,
            and 0.95 when a verb fired and an MCP annotation corroborated it;
      \item consult the model only below 0.5 confidence;
      \item ask the model where the rules lack semantic information, not where
            two rules disagree -- precedence resolves rule conflicts;
      \item every rules-decided tier remains attributable to the matched pattern
            and scorer;
      \item honest caveat: on the evaluated catalogs the fallback never fired,
            so \S\ref{sec:eval} must bound what can be claimed about it.
    \end{itemize}
\end{itemize}

% -------------------------------------------------------------------------
\subsubsection{Asset Sensitivity}
\label{sec:sensitivity}
% ~300 words.
% PURPOSE: explain how an asset receives a 1--5 sensitivity value even when the
% organization does not publish a numeric inventory.
\begin{itemize}
  \item \textbf{The problem.} A risk model needs to know how consequential loss
        or compromise of an asset would be, but organizations often do not
        maintain a complete, current, numeric sensitivity table for every
        resource.

  \item \textbf{Why a pre-scored inventory is a poor assumption.} It is costly to
        build, can itself expose sensitive information, and becomes stale as new
        resources are created.

  \item \textbf{What organizations do maintain.} A data/information
        classification policy with named classes, consequence language,
        recognition rules, and aggregation rules.

  \item \textbf{Classify then map.}
    \begin{itemize}
      \item the organization defines the classes and describes the consequence
            of loss in its own terms, without framework-specific numbers;
      \item \framework{} owns the 1--5 scale and maps the policy's adverse-impact
            language onto that scale;
      \item each asset/register row is then classified using its description and
            the organization's recognition and aggregation rules.
    \end{itemize}

  \item \textbf{Division of responsibility.} The organization says what counts
        as sensitive in its domain; the framework converts those classes into a
        common risk scale. Neither side is asked to perform the other's task.

  \item \textbf{Precedent.} FIPS~199 / SP~800-60 similarly begin from adverse
        impact and security categorization rather than requiring the risk model
        to invent business meaning from scratch.

  \item \textbf{Rule-vs.-row property.} Because classification is rule-based, the
        same policy can classify assets created after the policy was written.
        Present this as a property of the construction, not an evaluated result.

  \item \textbf{Evaluation separation.} Any per-asset ground-truth sensitivity
        table created for the study is an evaluation instrument and is not
        exposed to the scorer.
\end{itemize}
% OMIT: ISO/NIST/CIS comparison -- evaluation arm, not framework definition.

% -------------------------------------------------------------------------
\subsubsection{Blast Radius}
\label{sec:blast}
% ~300 words.
% PURPOSE: explain the scope/propagation factor for a (tool, asset) pairing.
\begin{itemize}
  \item \textbf{Unit and question.} Blast Radius is scored per (tool, asset)
        pairing: if this operation is exercised against this asset, how far can
        the consequences propagate?

  \item \textbf{The 1--5 rubric.} (1) almost nobody depends on it; (2) a few
        subjects; (3) a group attached to what was touched -- the normal case;
        (4) the whole population inside the asset boundary; (5) consequences
        can extend beyond that boundary.

  \item \textbf{Tier 5 needs a sanctioned propagation route} named by the asset
        register itself: for example, an authentication hub, a returned value
        usable elsewhere, one call reaching the entire population, or
        destruction with no practical restoration path.

  \item \textbf{Why not count items.} Item count is not consequence scope.
        Reading one credential file touches one item but may compromise every
        system the credential opens.

  \item \textbf{The floor and deliberate coupling.} Sensitivity and
        irreversibility may raise Blast Radius but never lower it. State clearly
        that this introduces one-directional dependence between primitives on
        purpose: a pinpoint item in a crown-jewel or irreversible path should
        not be treated as a harmless local effect.

  \item \textbf{The removed roof.} An earlier version capped blast in some cases;
        removing that roof avoids systematic under-scoring. Keep only if space
        allows.

  \item \textbf{The weak point.} Blast Radius is the noisiest of the three static
        factors. Do not hide this; \S\ref{sec:eval} should quantify its
        variability.
\end{itemize}

% -------------------------------------------------------------------------
\subsection{The Dynamic Layer}
\label{sec:dynamic}
% ~350-400 words.
% PURPOSE: explain what the static table cannot observe and how request-time
% evidence modifies the baseline risk assessment.
\begin{itemize}
  \item \textbf{The handover.} The static layer answers how consequential the
        capability is for this (tool, asset) pair. It cannot see the actual
        request, the agent's recent behavior, or the surrounding session. The
        dynamic layer adds those request-time facts.

  \item \textbf{Why dynamic evidence is needed.} Two invocations can share the
        same static cell yet differ materially: one may use ordinary arguments
        seen many times before, while another introduces a new target, a large
        scope, an unusual parameter value, or a suspicious sequence of actions.

  \item \textbf{Two factors and one bound.} The layer introduces no scale of
        its own; each factor moves a dimension the static product already has.
    \begin{itemize}
      \item \emph{likelihood} -- share of this deployment's traffic on this
            tool, read as surprisal and blended with a catalog prior
            ($\alpha = 30$) so there is no cold start; keyed
            (deployment, tool), clamped to $[0.1, 1.0]$; a tool never seen, a
            deployment too new to judge, or an unresolved resource all return
            1.0;
      \item \emph{escalation} -- three bins ($0.9 / 1.0 / 1.3$) from Tukey's
            fences over the argument values this deployment has passed;
            quantities only, never identifiers; below 20 observations it
            returns 1.0;
      \item \emph{floor} -- the lower boundary of the band the cell occupied
            before anything moved it, absolute rather than proportional, so
            the runtime reorders calls inside a band and may promote one, but
            never demotes a call out of its scanned band.
    \end{itemize}

  \item \textbf{Composition with the static baseline.} One equation, given as
        Eq.~\ref{eq:rho}: the static cell times likelihood times escalation,
        clamped below by the floor and above by 125. Only the static score can
        make the result large.

  \item \textbf{What is not in the layer.} The embedding likelihood, the
        behavioral-baseline and session-sequence band signals, and the optional
        LLM judge of earlier drafts are not part of this composition, and
        neither is a runtime egress test -- no predicate survived, and moving
        blast radius unconditionally manufactured false alarms. Session context
        is observable and unused. Cost stays negligible: no signal calls a
        model.

  \item \textbf{What the layer outputs.} The dynamic layer produces the final
        invocation-level assessment $\rho(c)$ together with the static factors
        and dynamic evidence that led to it, so the result remains auditable.

  \item \textbf{Scope boundary.} Stop here. The framework reports risk; it does
        not prescribe or evaluate an enforcement mechanism such as
        allow/review/deny thresholds.
\end{itemize}

% ---------------------------------------------------------------------------
% STRUCTURE NOTES
%
% * 4.1 should make the architecture obvious before introducing any rubric.
%   Its job is not just "what goes in / what comes out" but "why static and
%   dynamic are different kinds of evidence and therefore separate layers."
%
% * 4.2 is now the full static story. Keeping Tool Impact, Asset Sensitivity,
%   and Blast Radius as 4.2.1--4.2.3 mirrors Eq.~\ref{eq:score} and makes the
%   paper much easier to navigate.
%
% * The old Rules-to-Model Hand-off is part of Tool Impact because it only
%   explains how that factor is derived; it does not deserve a peer-level
%   subsection next to the three factors.
%
% * Enforcement is removed. If later sections discuss possible organizational
%   responses to a risk band, describe them as downstream uses/examples rather
%   than as part of the framework itself.
%
% * Consider one short worked example near the end of 4.1: one tool, one asset,
%   the three static factors, then one dynamic adjustment. Do NOT end the
%   example with allow/review/block because enforcement is outside scope.
% ---------------------------------------------------------------------------
